\documentclass[11pt]{article}
\usepackage[T1]{fontenc}
\usepackage{lmodern}
\usepackage[margin=1in]{geometry}
\usepackage{amsmath,amssymb,amsthm,mathtools}
\usepackage{microtype}
\usepackage[hidelinks]{hyperref}
\usepackage{cleveref}
\numberwithin{equation}{section}
\newtheorem{theorem}{Theorem}[section]
\newtheorem{proposition}{Proposition}[section]
\newtheorem{lemma}{Lemma}[section]
\theoremstyle{remark}
\newtheorem{remark}{Remark}[section]
\newcommand{\R}{\mathbb R}
\newcommand{\C}{\mathbb C}
\newcommand{\norm}[1]{\left\|#1\right\|}
\newcommand{\eps}{\epsilon}
\newcommand{\dH}{\dot H}
\newcommand{\Lam}[1]{\Lambda_{#1}}
\newcommand{\Pc}{P_c}
\newcommand{\Pt}[1]{\widetilde P_{#1}}
\DeclareMathOperator{\Rea}{Re}
\DeclareMathOperator{\Ima}{Im}
\DeclareMathOperator{\Ran}{Ran}
\DeclareMathOperator{\Lip}{Lip}
\title{Finite-time blowup for a Schr\"odinger envelope coupled to ionization}
\author{Kaiwen Jin\\\small\texttt{yui.ui.siki@gmail.com}}
\date{}
\begin{document}
\maketitle
\begin{abstract}
We construct a finite-time singularity for a three-dimensional focusing
Schr\"odinger envelope coupled to an ionization density. For one strictly
positive ionization coefficient, with cubic absorption and no avalanche
term, a nonzero compactly supported smooth optical datum and zero initial
density produce a smooth solution whose gradient norm tends to infinity
at a finite time. The density remains nonnegative. Thus global smooth
propagation for all admissible coefficients and Schwartz data is false.
The construction continues a scalar self-similar profile to positive
ionization, then solves the optical and density perturbations together.
A fixed scalar linear theory and compactly supported finite-dimensional
corrections allow the singular solution to start from compact smooth data.
\end{abstract}

\section{Introduction}\label{sec:introduction}

Ionization modifies a focusing nonlinear Schr\"odinger envelope through
both absorption and a density-dependent defocusing term. Such terms arise
in models of laser filamentation; their derivation and the question of
global regularization are discussed by Dumas, Lannes and
Szeftel~\cite[Section 4.1.2, equation (54) and Open Problem 3]{DLS}.
We consider the following precise system on $\R^3$, with spatial
coordinate $x=(x_1,x_2,z)$:
\begin{equation}\label{eq:original}
 \begin{split}
 i(\partial_t+c_g\partial_z)u+\eps\Delta u
          +\eps(|u|^2-\rho)u
   &=-i\eps c\bigl(\alpha_4|u|^{2K-2}u+\alpha_5\rho u\bigr),\\
 \partial_t\rho&=\eps\alpha_4|u|^{2K}
                       +\eps\alpha_5\rho|u|^2.
 \end{split}
\end{equation}
Here $u$ is complex, $\rho$ is real, the Laplacian is elliptic in all
three spatial coordinates, and
\[
 \eps,c,c_g,\alpha_4>0,\qquad \alpha_5\ge0,
 \qquad K\in\mathbb Z,\quad K\ge2.
\]
The universal global-propagation question asks whether every Schwartz
pair $(u_0,\rho_0)$ with $\rho_0\ge0$ has a solution in
$C([0,\infty);H^m\times H^m)$ for every integer $m\ge3$.
This is the formulation recorded as AIM Problem~441.\footnote{The
precise problem statement is available at
\href{https://github.com/MColbrook/AIM/blob/a153f469ee0ea64e876b2ad97c79d252c19a3e22/problems/441-ionization-laser-envelope-global-regularity.md}{the immutable AIM 441 statement}.}
The theorem below gives an admissible counterexample.

For the scalar focusing cubic NLS in three dimensions, Donninger and
Sch\"orkhuber~\cite{DS} construct a self-similar blowup profile by a
computer-assisted proof with exact fraction arithmetic. That profile
has finite energy but a spatial tail of order $r^{-1}$, hence is not in
$L^2$. Li~\cite{Li} develops linear dispersive estimates and finite
codimensional stability theory for a class of mass-supercritical
self-similar NLS profiles. We use the quantitative profile inverse from
the former work and the fixed-profile linear estimates and compactly
supported correction modes from the latter. Neither result directly
treats the ionization coupling in~\eqref{eq:original}.

\begin{theorem}\label{thm:main}
There exist $\delta>0$, a nonzero $u_0\in C_c^\infty(\R^3;\C)$, and
$T\in(0,\infty)$ such that, for
\[
 \eps=c=c_g=1,\qquad \alpha_4=\delta,\qquad K=2,
 \qquad \alpha_5=0,\qquad \rho_0=0,
\]
the original smooth solution of~\eqref{eq:original} exists on $[0,T)$,
belongs to $C_{\rm loc}([0,T);H^m(\R^3;\C)\times H^m(\R^3;\R))$
for every integer $m\ge3$, has $\rho(t,x)\ge0$, and satisfies
\[
 \lim_{t\uparrow T}\norm{\nabla u(t)}_{L^2(\R^3)}=\infty.
\]
In particular it has no global continuation in $C_t(H^3\times H^3)$.
\end{theorem}

One strictly positive coefficient and one Schwartz datum suffice to
negate the universal assertion. The theorem selects a finite member of
a sequence of small positive coefficients; it gives neither an explicit
full Cauchy threshold nor failure for every coefficient. It also makes
no open-set stability claim.

There are two obstacles to transferring the scalar construction. First,
the density is of the same similarity order as $|u|^2$, so it must be
included in the stationary profile. Its radial equation determines a
nonlocal integral of the fourth power of the optical profile. We bound
this added forcing in the exact spaces of the scalar inverse and preserve
the opposite signs of its real obstruction. A compactness argument then
selects one scalar limiting profile before any spectral constants are
chosen.

Second, a compact cutoff of the profile must be compatible with the
unstable finite coordinates, and the evolving density cannot be replaced
by its stationary value. We use modified projections whose finite ranges
consist of compactly supported smooth functions to make the optical datum
compact, and solve the
two perturbation equations jointly. Moving with the optical envelope
introduces a translation in the density equation; its characteristic
propagator retains that term exactly. A mixed space-time estimate controls
the quartic density source, and weighting the density perturbation by
the inverse square root of the ionization coefficient makes the feedback
contractive. The resulting solution approaches a nonzero profile in a
shrinking spatial scale, forcing its gradient norm to diverge.

Section~\ref{sec:profiles} constructs the profiles and verifies the
hypotheses of the scalar linear theory. Section~\ref{sec:estimates}
establishes the linear, product and transport estimates.
Section~\ref{sec:construction} constructs the compact-data solution and
identifies it with the smooth physical Cauchy solution.
Appendix~\ref{app:arithmetic} describes the finite exact calculation.
We write $\mathcal Ff(\xi)=(2\pi)^{-3/2}\int e^{-ix\cdot\xi}f(x)\,dx$.
For $0<s<3/2$, $\dH^s$ has norm
$\norm{|\xi|^s\mathcal Ff}_2$ and its canonical realization in
$L^{6/(3-2s)}$; $H^s$ uses the inhomogeneous Fourier weight.
All intersection norms are sums, and all spatial norms use Lebesgue
measure on $\R^3$.

\section{Profiles with positive ionization}\label{sec:profiles}

We construct profiles for the parameter choice in Theorem~\ref{thm:main}.
For a function of the similarity variable $y\in\R^3$, write
$\Lam a f=af+y\cdot\nabla f$.
The stationary equations to be solved first are
\begin{equation}\label{eq:stationary}
 \Delta Q-Q+i\beta\Lam1 Q+(1+i\delta)|Q|^2Q-RQ=0,
 \qquad \beta\Lam2 R=\delta|Q|^4.
\end{equation}
The density profile is determined by the optical profile. For a radial
$Q(y)=q(|y|)$, its regular solution is
\begin{equation}\label{eq:density-profile}
 R(r)=\frac{\delta}{\beta r^2}\int_0^r s|q(s)|^4\,ds.
\end{equation}
This identity makes the added coupling small when $\delta$ is small, while
keeping $\delta$ strictly positive.

\subsection{The quantitative scalar construction}\label{subsec:scalar-construction}

We use the concrete scalar construction of Donninger and
Sch\"orkhuber~\cite{DS}, with precisely the parameter and polynomial used
there. Set
\[
 \beta_* =\frac{772201763088846}{841768781900003},\qquad
 J=[-10^{-10},10^{-10}],\qquad r_*=1.2\cdot10^{-6},
 \qquad \beta=\beta_*+\theta\quad(\theta\in J).
\]
The radial change of variables is
\begin{equation}\label{eq:compactification}
 r=\frac{1+y}{1-y},\qquad
 q(r)=(1+r)^{-1-i/\beta}F(y),\qquad -1<y<1.
\end{equation}
In this subsection $y$ is the scalar compactified radial coordinate.
The source spaces, denoted here by $X_p,Y_p$, have norms
\[
 \norm f_{X_p}=\norm f_\infty+
       \sup_{-1<y<1}(1-y^2)|f'(y)|,
 \qquad
 \norm g_{Y_p}=\sup_{-1<y<1}(1+y)(1-y)^2|g(y)|.
\]
Functions in $X_p$ are continuous on $[-1,1]$, are $C^1$ in its interior,
and have continuously extended weighted derivative. The space $Y_p$ is
the source forcing space, which permits the fundamental-matrix join at
$y=0$. Continuous forcings with continuous weighted endpoint limits belong
to it.

Here is the precise part of the scalar construction that we need. In the
notation of~\cite{DS}, its $a_*,a$ become $\beta_*,\theta$. If
$P_*(y)=\sum_{n=0}^{50}c_nT_n(y)$ is the rational polynomial in its
Appendix~A, put
\begin{equation}\label{eq:fstar}
 f_{*,\theta}(y)=P_*(y)-P_*(1)
       +\frac{2i\beta}{1-i\beta}P_*'(1).
\end{equation}
There are a real-linear map $J_\theta:Y_p\to X_p$, a real functional
$\psi_\theta:Y_p\to\R$, an approximate differential operator
$\mathcal L_{\theta}$ induced by the source's chosen fundamental matrix,
and a nonlinear map $G_\theta$ with
\begin{equation}\label{eq:source-cancellation}
 G_\theta(f)=\mathcal L_{\theta}f-\mathcal R_\beta(f_{*,\theta}+f).
\end{equation}
Here $\mathcal R_\beta$ is the compactified residual of the scalar
equation obtained from~\eqref{eq:stationary} with $\delta=R=0$.
The second derivatives cancel in~\eqref{eq:source-cancellation}, so this
identity is meaningful for $f\in X_p$ before $C^2$ regularity is known.
The inverse property is the following: if
$f=J_\theta g$ and $\psi_\theta g=0$, then
$\mathcal L_{\theta}f=g$ away from the artificial join, and the source construction
gives $C^1$ matching across it. The scalar obstruction is exactly the term
removed in defining the regularized inverse.

The maps are strongly continuous in $\theta$ on fixed inputs. On the
closed $X_p$ ball of radius $r_*$, the quantitative estimates are
\begin{align}
 \norm{J_\theta}_{Y_p\to X_p}&\le232,
 &\norm{\psi_\theta}_{Y_p\to\R}&\le13,\label{eq:source-inverse}\\
 \norm{G_\theta(f)}_{Y_p}
 &\le5\cdot10^{-6}\norm f+10\norm f^2+3\norm f^3+5\cdot10^{-9},
 \label{eq:source-nonlinear}\\
 \norm{G_\theta(f)-G_\theta(g)}_{Y_p}
 &\le\bigl[5\cdot10^{-6}+10(\norm f+\norm g)
       +3(\norm f^2+\norm g^2)\bigr]\norm{f-g}.
 \label{eq:source-lipschitz}
\end{align}
Unsubscripted correction norms here are $X_p$ norms. At
$\theta_-=-10^{-10}$ and $\theta_+=10^{-10}$, respectively,
\begin{equation}\label{eq:source-signs}
 \psi_{\theta_-}\mathcal R_{\beta_*+\theta_-}(f_{*,\theta_-})
 \ge4\cdot10^{-10},\qquad
 \psi_{\theta_+}\mathcal R_{\beta_*+\theta_+}(f_{*,\theta_+})
 \le-4\cdot10^{-10}.
\end{equation}
These are results for the scalar operator, rather than a coupled profile
theorem. The inverse and continuity statements are Definition~3.1 and
Lemmas~3.27, 3.29, 3.35, 3.37--3.39 of~\cite{DS}; the bounds are
Lemma~7.5, Proposition~7.6, Lemma~10.1 and Proposition~10.2. The strict
signs are Corollary~9.23. We use the weaker residual allowance displayed
in Proposition~10.2; Proposition~6.4 gives $4\cdot10^{-9}$.

\subsection{Adding the density and absorption terms}

For $F=f_{*,\theta}+f$, define $q$ by~\eqref{eq:compactification} and
$R_F$ by~\eqref{eq:density-profile}. Since $dy/dr=(1-y)^2/2$,
the additional compactified residual is exactly
\begin{equation}\label{eq:added-residual}
 P_{\delta,\theta}(F)
   =\frac{i\delta F|F|^2}{(1-y)^2}
       -\frac{4R_F(r(y))F}{(1-y)^4}.
\end{equation}
Indeed, dividing the additional physical terms by
$(1+r)^{-1-i/\beta}(dy/dr)^2$ gives this formula.
If $M\ge\norm F_\infty$ and $\beta_{\min}=\beta_*-10^{-10}>9/10$,
then
\begin{align}
 0\le R_F(r)&\le\frac{2\delta M^4}{\beta(1+r)^2},\label{eq:R-bound}\\
 \norm{P_{\delta,\theta}(F)}_{Y_p}
   &\le\delta(2M^3+4M^5/\beta_{\min}),\label{eq:P-bound}\\
 \norm{P_{\delta,\theta}(F)-P_{\delta,\theta}(G)}_{Y_p}
   &\le\delta(6M^2+20M^4/\beta_{\min})\norm{F-G}_{X_p}.
   \label{eq:P-difference}
\end{align}
To see the density bound, put $I(r)=\int_0^r s(1+s)^{-4}\,ds$.
For $r\le1$, $I(r)\le r^2/2$; for $r\ge1$, $I(r)\le1/6$.
Thus $(1+r)^2r^{-2}I(r)\le2$. Multiplication by the $Y_p$ weight
in~\eqref{eq:added-residual} gives~\eqref{eq:P-bound}; for the density
term the remaining factor is $4(1+y)/(1-y)^2=2r(1+r)$.
Factoring the cubic and fourth-power differences gives
\eqref{eq:P-difference}.
At $r=0$ the weighted density term tends to zero. At infinity the integral
$\int_0^r s|q(s)|^4\,ds$ converges absolutely, so the weighted forcing has
a finite endpoint limit. Thus the added residual belongs to $Y_p$.

The finite rational inequalities proved in Appendix~\ref{app:arithmetic}
give, throughout the correction ball,
\begin{equation}\label{eq:profile-envelopes}
 \norm F_\infty<50,\qquad \Rea F(-1)<-1,
 \qquad |f_{*,\theta}(1)|>1/10.
\end{equation}
For $0<\delta\le10^{-23}$, equations~\eqref{eq:P-bound} and
\eqref{eq:P-difference} therefore give the convenient bounds
$p=10^{-13}$ and $l_P=10^{-6}$ for the forcing and its Lipschitz constant.
Define
\[
 T_{\delta,\theta}f
  =J_\theta\bigl[G_\theta(f)
       -P_{\delta,\theta}(f_{*,\theta}+f)\bigr].
\]
The self-map bound is
\[
 232(5\cdot10^{-6}r_*+10r_*^2+3r_*^3+5\cdot10^{-9}+p)<r_*.
\]
Its Lipschitz constant is at most
\begin{equation}\label{eq:profile-contraction}
 232(5\cdot10^{-6}+20r_*+6r_*^2+l_P)<1/100.
\end{equation}
The contraction theorem gives a unique $f_\theta^\delta$ in this ball,
continuous in $\theta$ and $\delta$, including $\delta=0$.
Here strong continuity of $J_\theta$ on fixed inputs suffices: subtract
the two fixed-point equations and absorb their common contraction constant.

For $\delta=0$, equations~\eqref{eq:source-nonlinear}--\eqref{eq:source-signs}
show that $\psi_\theta G_\theta(f_\theta^0)$ has negative and positive
endpoint signs, with margin
\[
 m_0=4\cdot10^{-10}
       -13(5\cdot10^{-6}r_*+10r_*^2+3r_*^3)>10^{-10}.
\]
Moreover, $\norm{f_\theta^\delta-f_\theta^0}_{X_p}<235p$.
Consequently the continuous obstruction
\[
 \Psi_\delta(\theta)=\psi_\theta\bigl[G_\theta(f_\theta^\delta)
       -P_{\delta,\theta}(f_{*,\theta}+f_\theta^\delta)\bigr]
\]
changes at either endpoint by at most
\[
 13\left[1+
 \frac{232(5\cdot10^{-6}+20r_*+6r_*^2)}{1-1/100}\right]p<10^{-11}.
\]
Its signs persist. The intermediate value theorem gives a root
$\theta_\delta\in J$ for each such positive $\delta$.

At the root the regularization term vanishes. The inverse property and
\eqref{eq:source-cancellation} imply
$\mathcal R_\beta(F)+P_{\delta,\theta}(F)=0$ on the two sides of the join.
The correction is $C^1$ there. After cancellation of the artificial
coefficients, its physical second-order ODE has continuous forcing across
the join, and hence gives a $C^2$ solution. We obtain
\begin{equation}\label{eq:radial-profile}
 q''+(2/r+i\beta r)q'+(i\beta-1)q
          +(1+i\delta)|q|^2q-R_Fq=0.
\end{equation}
This establishes the coupled profile by the scalar inverse and a new
forcing estimate; no differentiability of the selected parameter root is
needed.

\subsection{Regularity and a fixed scalar limit}\label{subsec:profile-limit}

The source norm implies uniformly
$|q(r)|\le C(1+r)^{-1}$ and $|q'(r)|\le C/[r(1+r)]$.
At the origin put
$A(q,R)=-(i\beta-1)q-(1+i\delta)|q|^2q+Rq$.
Integrating~\eqref{eq:radial-profile} gives
\begin{align}
 q'(r)&=r\int_0^1\tau^2 e^{i\beta r^2(\tau^2-1)/2}
            A(q(r\tau),R(r\tau))\,d\tau,\label{eq:origin-q}\\
 R(x)&=\frac{\delta}{\beta}\int_0^1\tau|Q(\tau x)|^4\,d\tau.
 \label{eq:origin-R}
\end{align}
The possible integration constant in
$(r^2e^{i\beta r^2/2}q')'=r^2e^{i\beta r^2/2}A$ vanishes because
$q'=O(r^{-1})$. In Cartesian coordinates~\eqref{eq:origin-q} says
that $\nabla Q(x)$ is $x$ times its continuous integral. It first gives
$Q\in C^1$ and $q'=O(r)$ at zero. If $Q,R$ are $C^k$ locally, these
compact-parameter formulas give $Q\in C^{k+1}$ and then $R\in C^{k+1}$.
Differentiation is uniform on compact products. Induction proves smoothness
across the origin, with uniform bounds at each fixed local derivative order.
At infinity, differentiating~\eqref{eq:density-profile} twice gives
\begin{equation}\label{eq:R-tail}
 R=O(\delta r^{-2}),\quad R'=O(\delta r^{-3}),\quad
 R''=O(\delta r^{-4}),\qquad
 \norm R_{H^2}+\norm R_\infty\le C\delta.
\end{equation}

Choose $\delta_j\downarrow0$ and one root $\theta_j$ for each.
After a subsequence $\theta_j\to\theta_0$. Uniform contraction continuity
gives convergence in $X_p$ to $f_{\theta_0}^0$, and the limiting obstruction
is zero. Let $\beta_0=\beta_*+\theta_0$ and let $Q_0$ be this scalar
profile. The two endpoint bounds~\eqref{eq:profile-envelopes} make it
nontrivial and give a nonzero coefficient at infinity.

We fix $\beta_0$ for the remainder of the proof. Rescale the positive
profiles, initially denoted $q_j,R_j^{\rm old}$, by
\[
 d_j=(\beta_0/\beta_j)^{1/2},\qquad
 Q_j(y)=d_jq_j(d_j|y|),\qquad
 R_j(y)=d_j^2R_j^{\rm old}(d_j|y|),\qquad \omega_j=d_j^2.
\]
Their equations become
\begin{equation}\label{eq:rescaled-profiles}
 \Delta Q_j-\omega_jQ_j+i\beta_0\Lam1 Q_j
 +(1+i\delta_j)|Q_j|^2Q_j-R_jQ_j=0,
 \qquad \beta_0\Lam2 R_j=\delta_j|Q_j|^4.
\end{equation}
For the first equation every residual term acquires $d_j^3$.
For the second,
$\beta_0\Lam2 R_j=\beta_jd_j^4(\Lam2 R_j^{\rm old})(d_jy)$;
the fourth power on the right acquires the same factor.
We have $\omega_j\to1$ and~\eqref{eq:R-tail} with $\delta=\delta_j$.
Local convergence of profiles and first derivatives, followed by splitting
off their uniform tails, gives
\begin{equation}\label{eq:coefficient-convergence}
 \begin{split}
 &\norm{|Q_j|^2-|Q_0|^2}_\infty+
 \norm{|Q_j|^2-|Q_0|^2}_3+
 \norm{\nabla(|Q_j|^2-|Q_0|^2)}_3\longrightarrow0,\\
 &\norm{Q_j^2-Q_0^2}_\infty+
 \norm{Q_j^2-Q_0^2}_3+
 \norm{\nabla(Q_j^2-Q_0^2)}_3\longrightarrow0.
 \end{split}
\end{equation}
The tail powers are $r^{-2}$ and $r^{-3}$, both cubically integrable in
three dimensions. Uniformly, $Q_j\in L^p$ for every $p>3$ and
$\nabla Q_j\in L^2\cap L^3\cap L^\infty$.

\subsection{The scalar profile hypotheses for the linear estimates}
\label{subsec:scalar-tail}

Li's linear theory~\cite{Li} requires a smooth nowhere-zero scalar
profile, nonzero limiting modulus of $rQ_0$, and symbol bounds for its
derivatives. We verify these properties for the particular fixed limit.
Write $Q_0(x)=q(|x|)$, $\beta=\beta_0$,
$\kappa=1+i/\beta$ and $g=r^\kappa q$. Endpoint continuity gives
$g\to g_\infty\ne0$, while the first derivative bound gives $g'=O(r^{-1})$.
The scalar equation becomes
\[
 g''+Dg'=-r^{-2}B(g),\qquad
 D=i\beta r-\frac{2i}{\beta r},\qquad
 B(g)=\kappa(\kappa-1)g+|g|^2g.
\]
For $\mathfrak m=e^{i\beta r^2/2}r^{-2i/\beta}$ we have
$\mathfrak m'/\mathfrak m=D$ and $|\mathfrak m|=1$. Since
$(\mathfrak m g')'$ is absolutely integrable and $g'\to0$,
\begin{equation}\label{eq:tail-integral}
 g'(r)=\mathfrak m(r)^{-1}\int_r^\infty
                  \mathfrak m(s)s^{-2}B(g(s))\,ds.
\end{equation}
For large $r$, $|D|\ge cr$. Put $f=r^{-2}B(g)$; initially
$f=O(r^{-2})$, $f'=O(r^{-3})$. Integration by parts using
$\mathfrak m'=D\mathfrak m$ gives $g'=O(r^{-3})$.
The equation then gives $g''=O(r^{-2})$, and hence $f''=O(r^{-4})$.
A second integration by parts of the remainder gives
\[
 g'=-f/D+O(r^{-5}),\qquad g''=O(r^{-4}),\qquad
 g-g_\infty=O(r^{-2}).
\]
The boundary terms vanish, since $f/D=O(r^{-3})$ and its next
derivative is $O(r^{-4})$.

To obtain all derivative orders without differentiating an uncontrolled
remainder, let $S_m^1(R)$ be the complete weighted $C^1$ space with norm
\[
 \norm f_{S_m^1}=\sup_{r\ge R}
            (r^m|f(r)|+r^{m+1}|f'(r)|),\qquad m>0.
\]
The integral operator
\[
 \mathcal V f=-\int_r^\infty\mathfrak m(t)^{-1}
                 \int_t^\infty\mathfrak m(s)f(s)\,ds\,dt
\]
is defined with the inner oscillatory integral first. Integration by parts
bounds that integral by $C_mr^{-m-1}\norm f_{S_m^1}$; the outer integral
is absolutely convergent. Thus $\mathcal V$ is bounded on $S_m^1$ and
$(\mathcal V f)''+D(\mathcal V f)'=-f$.
For a bounded real-linear coefficient $A(r)$ with $A'=O(r^{-1})$,
\begin{equation}\label{eq:volterra-small}
 \norm{\mathcal V(r^{-2}Af)}_{S_m^1}
             \le C_m R^{-2}\norm f_{S_m^1}.
\end{equation}
Completeness follows by uniform convergence of weighted values and
derivatives and the fundamental theorem of calculus on compact intervals.

For each integer $N$, choose
$G_N=g_\infty+\sum_{\ell=1}^Nc_\ell r^{-2\ell}$ to cancel the residual
$G_N''+DG_N'+r^{-2}B(G_N)$ through order $r^{-2N}$.
At the new order the coefficient of $c_\ell$ is $-2\ell i\beta\ne0$.
The residual lies in $S_{2N+2}^1$. The remainder equation is
\[
 h=\mathcal V\{r^{-2}[B(G_N+h)-B(G_N)]
                  +G_N''+DG_N'+r^{-2}B(G_N)\}.
\]
By~\eqref{eq:volterra-small} it is a contraction in a ball of
$S_{2N+2}^1$ for sufficiently large $R$. The actual remainder $g-G_N$
already lies in $S_2^1$. The difference $v$ between it and the constructed
remainder solves $v=\mathcal V(r^{-2}Av)$ on $S_2^1$.
Here $A$ is the real derivative of the polynomial $B$ averaged along the
segment between the two solutions. Their bounded values and derivatives
give bounded $A,rA'$. Enlarging $R$ in~\eqref{eq:volterra-small} forces
$v=0$. Therefore $g-G_N\in S_{2N+2}^1$ for every $N$.
Differentiation of its smooth ODE gives, for $k\ge2$,
$h^{(k)}=O(r^{-2N-4+k})$. For any specified derivative order choose $N$
sufficiently large and differentiate $q=r^{-\kappa}(G_N+h)$. This proves
\begin{equation}\label{eq:all-tail}
 |q^{(k)}(r)|\le C_k r^{-1-k}\quad(r\ge1, k\ge0).
\end{equation}
The oscillatory homogeneous branch is excluded already by $g'\to0$.

Finally let $n=|q|^2$ and $W=\Ima(q'\bar q)$. The imaginary part of the
radial equation, integrated from zero, is
\[
 r^2W(r)=-\frac{\beta}{2}r^3n(r)
                +\frac{\beta}{2}\int_0^r s^2n(s)\,ds.
\]
A positive-radius zero would make the nonnegative integral vanish, which
contradicts the nonzero origin value in~\eqref{eq:profile-envelopes}.
Thus $q$ never vanishes and $\lim r|q(r)|=|g_\infty|>0$.
Together with smoothness and~\eqref{eq:all-tail}, this verifies
conditions (1.5)--(1.8) of~\cite{Li} for $d=p=3$ and $s_c=1/2$,
including the higher-derivative condition as printed there. Here $d$ is
the spatial dimension, $p$ is the exponent in $|u|^{p-1}u$, and
$s_c=d/2-2/(p-1)$.

\section{Estimates in similarity variables}\label{sec:estimates}

\subsection{A Green operator for one fixed scalar profile}

From now on $y\in\R^3$ is the similarity variable. All homogeneous
Sobolev spaces use the canonical realization given by Sobolev embedding;
intersection norms are sums. For the single fixed profile $Q_0,\beta_0$
from Section~\ref{sec:profiles}, let $\delta_1>0$ be the spectral gap
in Li's decomposition~\cite[Lemma 6.3]{Li}. Choose
\begin{equation}\label{eq:sigma-nu}
 \frac12<\sigma<\min\left(0.51,\frac12+\delta_1\right),
 \quad X=\dH^\sigma\cap\dH^1,\quad
 \mu=\beta_0(\sigma-1/2),\quad
 0<\nu<\min(\mu,\beta_0/2).
\end{equation}
All these choices precede the final choice of $j$.
The real-linear scalar operator is
\[
 L_0e=\Delta e-e+i\beta_0\Lam1 e+2|Q_0|^2e+Q_0^2\bar e.
\]
For conjugate pairs its matrix is
\begin{equation}\label{eq:H0}
 H_0=
 \begin{pmatrix}\Delta_{\beta_0}-1&0\\0&-\Delta_{-\beta_0}+1\end{pmatrix}
 -\frac{i\beta_0}{2}
 +\begin{pmatrix}2|Q_0|^2&Q_0^2\\-\bar Q_0^2&-2|Q_0|^2\end{pmatrix},
 \qquad \Delta_b=\Delta+ib(3/2+y\cdot\nabla).
\end{equation}
This is Li's operator (1.22)--(1.23) at $d=p=3$.
Its unbounded generator and flow are those in that theory.
Pair spaces have component-sum norms. The involution
$\mathcal J(z_1,z_2)=(\bar z_2,\bar z_1)$ has fixed pairs
$\iota e=(e,\bar e)$, of norm $2\norm e_X$.

We state the linear facts used below. Li's Lemmas~6.3, 7.1 and 7.2 give
original spectral projections $P_c,P_s,P_u$ and compactly modified
projections $\Pt c,\Pt s,\Pt u$, commuting with $\mathcal J$, such that
\begin{align}
 \Pt\alpha P_\alpha&=\Pt\alpha,&
 P_\alpha\Pt\alpha&=P_\alpha\quad(\alpha=s,u),\label{eq:finite-projections}\\
 \Pt cP_c&=P_c,&P_c\Pt c&=\Pt c.
 \label{eq:central-projections}
\end{align}
The modified finite projections are bounded on $X$ with finite-dimensional
ranges in $C_c^\infty$. Also $\Pt c=I-\Pt s-\Pt u$ is bounded on $X$,
and $\Ran\Pt c=\Ran P_c$.
When a projection acts on a scalar, it means its first component after
application to $\iota e$; commutation with $\mathcal J$ makes this a
real projection. The fixed factor two is absorbed in constants.
For $z\in X\cap\Ran P_c$ the homogeneous estimate is
\begin{equation}\label{eq:essential-flow}
 \norm{e^{isH_0}z}_{C([k,k+1];X)}
 +\norm{\nabla e^{isH_0}z}_{L^3([k,k+1];L^{18/5})}
 \le C e^{-\mu k}\norm z_X,
 \qquad k\in\mathbb Z_{\ge0}.
\end{equation}
Indeed, (6.2) of~\cite{Li} allows $18/5$ when $\sigma<0.51$;
$(3,18/5)$ is Schr\"odinger admissible, and the exponential weight in
its Lemma~7.2 gives~\eqref{eq:essential-flow} on each unit interval.
Time translation gives the same estimate for a later initial time.

Define the solution and forcing norms
\begin{align}
 \norm e_{\mathcal S_\nu}
 &=\sup_{k\ge0}e^{\nu k}
   \left(\norm e_{C([k,k+1];X)}+
          \norm{\nabla e}_{L^3([k,k+1];L^{18/5})}\right),
 \label{eq:S-norm}\\
 \norm F_{\mathcal F_\nu}
 &=\sup_{k\ge0}e^{\nu k}\norm F_{L^1([k,k+1];X)}.
 \label{eq:F-norm}
\end{align}
The central-stable initial space is $(I-\Pt u)X$.

\begin{lemma}\label{lem:green}
There is a bounded real-linear solution operator $\mathcal B$ for
$ie_s+L_0e=F$, with prescribed central-stable initial component
$e_{cs,0}\in(I-\Pt u)X$, such that
\begin{equation}\label{eq:green-bound}
 \norm{\mathcal B(e_{cs,0},F)}_{\mathcal S_\nu}
       \le C_B(\norm{e_{cs,0}}_X+\norm F_{\mathcal F_\nu}).
\end{equation}
Its remaining initial component lies in $\Ran\Pt u\subset C_c^\infty$.
All constants and projections depend on $Q_0,\beta_0,\sigma,\nu$, not on $j$.
\end{lemma}

\begin{proof}
Choose real bases $(\widetilde\varphi_{\alpha,\ell})_\ell$ of the
$\mathcal J$-fixed parts of $\Ran\Pt\alpha$ for $\alpha=s,u$.
The finite coordinate vectors $z_\alpha$ reconstruct the pair-valued modes
\[
 \widetilde Z_\alpha(s)=\sum_\ell(z_\alpha(s))_\ell
                   \widetilde\varphi_{\alpha,\ell},\qquad \alpha=s,u.
\]
The central component is $Z_c=\Pt c\iota e$.
In Li's modified coordinate equations (7.14), the finite coordinates obey
$z_\alpha'+A_\alpha z_\alpha=f_\alpha$.
Here $A_\alpha$ realizes $-i\Pt\alpha H_0P_\alpha$ in real finite
coordinates, and $f_\alpha$ realizes $-i\Pt\alpha(F,-\bar F)$.
The stable eigenvalues have real part at least $\beta_0\delta_1>\mu$,
whereas the unstable eigenvalues have real part at most zero. Use the
forward solution for the stable equation and
\[
 z_u(s)=-\int_s^\infty e^{A_u(\tau-s)}f_u(\tau)\,d\tau
\]
for the unstable one. On a Jordan block the exponential is an exponential
times a polynomial. The stable forward kernel is bounded by
$Ce^{-\gamma_sd}$ for some $\gamma_s>\mu$; the backward unstable kernel
grows at most polynomially. Their convolutions with $e^{-\nu k}$ obey
the required weighted bounds.

The central forcing is $\Pt c(F,-\bar F)$ together with
$-\Pt cH_0(\widetilde Z_s+\widetilde Z_u)$. Every finite modified basis
vector is compact and smooth, so its image under $H_0$ and then $\Pt c$
belongs to $X$. Thus this additional forcing is controlled by the finite
coordinates. Apply~\eqref{eq:essential-flow} to each forcing impulse,
use Minkowski's inequality, and sum
\begin{equation}\label{eq:exponential-sum}
 \sum_{\ell=0}^k e^{-\mu(k-\ell-1)_+}e^{-\nu\ell}
                   \le C_{\mu,\nu}e^{-\nu k}.
\end{equation}
An impulse inside the observation interval uses the unweighted homogeneous
estimate on the remaining subinterval. Its norm is bounded by that
interval's $L^1X$ forcing norm. The finite-coordinate sums are treated
with the same geometric sum, or a convergent polynomial-exponential sum
in the backward direction. This proves~\eqref{eq:green-bound}.

The scalar real structure and signs are consistent: direct calculation
gives $\mathcal JH_0\mathcal J=-H_0$, so $e^{isH_0}$ commutes with
$\mathcal J$. The pair forcing $(F,-\bar F)$ is anti-invariant, and
its generator forcing $-i(F,-\bar F)$ is invariant. The modified
projections preserve this fixed real subspace. The construction therefore
returns an actual scalar $e$ and its conjugate.
\end{proof}

The forcing norm in~\eqref{eq:F-norm} permits a small constant multiple
of $e$, since each interval has length one. This will accommodate the
frequency change $\omega_j-1$ without requiring an $L^2$ bound for the
similarity perturbation.

\subsection{Products at the quartic density nonlinearity}

For $I=[k,k+1]$ abbreviate
\[
 a_I(e)=\norm e_{C(I;X)}+
               \norm{\nabla e}_{L^3(I;L^{18/5})}.
\]
Sobolev embedding and interpolation give $X\hookrightarrow L^4\cap L^6$.
We also need an estimate combining the derivative Strichartz norm and
$L^6$. Averaging a smooth $f$ on a radius-$R$ ball and integrating its
gradient along segments gives
\[
 |f(x)|\le C R^{-1/2}\norm f_6+
       C\int_{|x-z|<R}|x-z|^{-2}|\nabla f(z)|\,dz
 \le C\bigl(R^{-1/2}\norm f_6+R^{1/6}\norm{\nabla f}_{18/5}\bigr).
\]
The kernel belongs locally to $L^{18/13}$, since $36/13<3$.
Balancing $R=(\norm f_6/\norm{\nabla f}_{18/5})^{3/2}$ gives
\begin{equation}\label{eq:mixed-GN}
 \norm f_\infty\le C\norm{\nabla f}_{18/5}^{3/4}\norm f_6^{1/4}.
\end{equation}
The zero cases follow by a limit. Mollification and an almost-everywhere
subsequential limit in $L^6$ extend the estimate to its stated norm class.
Interpolation also gives
$\norm{\nabla e}_3\le\norm{\nabla e}_2^{1/4}
\norm{\nabla e}_{18/5}^{3/4}$.
Raising to the fourth power and integrating in time yields
\begin{equation}\label{eq:unit-mixed}
 \norm e_{L^4(I;L^\infty)}+
 \norm{\nabla e}_{L^4(I;L^3)}+
 \norm e_{C(I;L^4\cap L^6)}\le C a_I(e).
\end{equation}

Recall from the end of Section~\ref{subsec:profile-limit} that $Q_j$
is uniformly bounded, belongs uniformly to $L^p$ for each fixed $p>3$,
and has gradient uniformly in $L^2\cap L^3\cap L^\infty$. Set
\[
 S_j(e)=|Q_j+e|^4-|Q_j|^4,\qquad
 N_j(e)=(1+i\delta_j)(2Q_j|e|^2+\bar Q_je^2+|e|^2e).
\]
Writing $a=a_I(e)$ and $b=\norm h_{C(I;H^1)}$, we have
\begin{align}
 \norm{S_j(e)}_{L^1(I;H^1)}&\le C(a+a^2+a^3+a^4),\label{eq:quartic}\\
 \norm{N_j(e)}_{L^1(I;H^1)}&\le C(a^2+a^3),\label{eq:cubic}\\
 \norm{h(Q_j+e)}_{L^1(I;H^1)}&\le Cb(1+a).
 \label{eq:density-product}
\end{align}
Here and below these constants are uniform for all sufficiently large $j$.
For clarity, expand
$|Q+e|^2=|Q|^2+2\Rea(Q\bar e)+|e|^2$ and square.
The linear term of $S_j$ is in $L^2$ by $Q_j^3\in L^3$ and $e\in L^6$.
Its derivative uses $Q_j^3\nabla e$ and $Q_j^2\nabla Q_j\,e$.
The quadratic undifferentiated terms use $e\in L^4$, and their derivatives
use $\norm e_\infty\norm{\nabla e}_2$.
The cubic terms use $\norm e_6^3$ and
$\norm e_\infty^2\norm{\nabla e}_2$.
At degree four,
\[
 \norm{e^4}_2\le\norm e_6^3\norm e_\infty,
 \qquad
 \norm{\nabla(e^4)}_2\le4\norm e_\infty^3\norm{\nabla e}_2.
\]
Derivatives on profile factors are controlled by their uniform bounds.
Time H\"older and~\eqref{eq:unit-mixed} prove~\eqref{eq:quartic}; the
quadratic and cubic terms also prove~\eqref{eq:cubic}. For the density
product use the bounded profile and
\[
 \norm{he}_2\le\norm h_3\norm e_6,\quad
 \norm{h\nabla e}_2\le\norm h_6\norm{\nabla e}_3,\quad
 \norm{e\nabla h}_2\le\norm e_\infty\norm{\nabla h}_2.
\]
These $H^1$ forcings belong to $X$, because $0<\sigma<1$.

The estimates also hold for differences. Factoring a difference of
products as a sum with one difference factor, before and after one
derivative, gives
\begin{align}
 \norm{S_j(e_1)-S_j(e_2)}_{L^1H^1}
  &\le C(1+a_1+a_2)^3 a_I(e_1-e_2),\label{eq:quartic-difference}\\
 \norm{N_j(e_1)-N_j(e_2)}_{L^1X}
  &\le C\bigl(a_1+a_2+(a_1+a_2)^2\bigr)a_I(e_1-e_2),
  \label{eq:cubic-difference}\\
 \norm{h_1(Q_j+e_1)-h_2(Q_j+e_2)}_{L^1X}
  &\le C\bigl[\norm{h_1-h_2}_{CH^1}(1+a_1)
                  +\norm{h_2}_{CH^1}a_I(e_1-e_2)\bigr].
  \label{eq:density-difference}
\end{align}
The omitted time interval in these three norms is $I$, and $a_i=a_I(e_i)$.
Finally, for a coefficient $V$,
\begin{equation}\label{eq:multiplier}
 \norm{Ve}_{H^1}
 \le C(\norm V_3+\norm V_\infty+\norm{\nabla V}_3)\norm e_X.
\end{equation}
Equations~\eqref{eq:coefficient-convergence}, \eqref{eq:R-tail} and
\eqref{eq:multiplier} make the profile-potential changes small from $X$
to $H^1$. The constant multiplier $\omega_j-1$ is instead small on $X$
itself. This distinction is useful because $Q_j$ and the initial cutoff
correction need not be in $L^2$ separately.

\subsection{The density propagator in the optical frame}

Let $e_z=(0,0,1)$ and introduce the optical frame
$X_{\rm ph}=x-te_z$,
$\widehat u(t,X_{\rm ph})=u(t,X_{\rm ph}+te_z)$, with the same convention
for $\widehat\rho$. The specialized original equations become
\[
 i\widehat u_t+\Delta\widehat u
       +(1+i\delta_j)|\widehat u|^2\widehat u
       -\widehat\rho\widehat u=0,
 \qquad
 \widehat\rho_t=\delta_j|\widehat u|^4+\partial_z\widehat\rho.
\]
For $0<\lambda_0\le1$ define
\begin{equation}\label{eq:similarity}
 \lambda(s)=\lambda_0e^{-\beta_0s},\quad
 t(s)=\frac{\lambda_0^2}{2\beta_0}(1-e^{-2\beta_0s}),\quad
 y=X_{\rm ph}/\lambda(s),\quad T=\frac{\lambda_0^2}{2\beta_0}.
\end{equation}
Since $dt/ds=\lambda^2$, the ansatz
\begin{equation}\label{eq:ansatz}
 \widehat u=\lambda^{-1}e^{i\omega_js}(Q_j+e),\qquad
 \widehat\rho=\lambda^{-2}(R_j+h)
\end{equation}
and~\eqref{eq:rescaled-profiles} give the exact density equation
\begin{equation}\label{eq:h-equation}
 h_s+\beta_0\Lam2 h-\lambda(s)\partial_zh
     =\delta_j S_j(e)+\lambda(s)\partial_zR_j,
 \qquad h(0)=-R_j.
\end{equation}
Thus the optical translation does not disappear from the density equation.
On compact intervals before $T$, the coordinate change is a smooth
diffeomorphism with nonzero Jacobian; the calculation also holds
distributionally by pulling back test functions.

Solving the characteristic equation $y'=\beta_0y-\lambda(s)e_z$ gives,
for $s\ge\tau$ and $d=s-\tau$,
\begin{equation}\label{eq:density-flow}
 (\mathcal T(s,\tau)f)(y)=e^{-2\beta_0d}
  f\left(e^{-\beta_0d}y+
       \frac{\lambda(\tau)}{2\beta_0}(1-e^{-2\beta_0d})e_z\right).
\end{equation}
Change of variables gives
\begin{equation}\label{eq:density-decay}
 \norm{\mathcal Tf}_2=e^{-\beta_0d/2}\norm f_2,
 \qquad
 \norm{\nabla\mathcal Tf}_2=e^{-3\beta_0d/2}\norm{\nabla f}_2.
\end{equation}
The propagator is strongly continuous on $H^1$, first on Schwartz
functions and then by density and these norm bounds.
Let $\mathcal Y_\nu$ be the real space $C([0,\infty);H^1)$ with norm
$\norm h_{\mathcal Y_\nu}=\sup_{s\ge0}e^{\nu s}\norm{h(s)}_{H^1}$.
Denote the integral solution of~\eqref{eq:h-equation} by $\Phi_\rho(e)$.
Equations~\eqref{eq:quartic}, \eqref{eq:density-decay} and
$\nu<\beta_0/2$ imply
\begin{align}
 \norm{\Phi_\rho(e)}_{\mathcal Y_\nu}
 &\le C\delta_j(1+\lambda_0+E+E^2+E^3+E^4),\label{eq:rho-map}\\
 \norm{\Phi_\rho(e_1)-\Phi_\rho(e_2)}_{\mathcal Y_\nu}
 &\le C\delta_j(1+E_1+E_2)^3\norm{e_1-e_2}_{\mathcal S_\nu},
 \label{eq:rho-map-difference}
\end{align}
where $E=\norm e_{\mathcal S_\nu}$ and $E_i=\norm{e_i}_{\mathcal S_\nu}$.
For example the background forcing is bounded in $H^1$ by
$C\delta_j\lambda_0e^{-\beta_0s}$, since
$\norm{\partial_zR_j}_{H^1}\le C\delta_j$.
The other terms are summed interval by interval using
\eqref{eq:exponential-sum} with density decay rate $\beta_0/2$ in place
of $\mu$. The map preserves real-valued $h$.
\section{A singular solution from compact smooth data}\label{sec:construction}

\subsection{Solving the coupled perturbation equations}

With the profiles and spaces now fixed, put
\[
 L_je=\Delta e-\omega_je+i\beta_0\Lam1 e
 +(2|Q_j|^2-R_j)e+Q_j^2\bar e
       +i\delta_j(2|Q_j|^2e+Q_j^2\bar e).
\]
Subtracting~\eqref{eq:rescaled-profiles} from the optical equation in
\eqref{eq:ansatz} gives
\begin{equation}\label{eq:e-equation}
 ie_s+L_0e=-(L_j-L_0)e-N_j(e)+h(Q_j+e)=:F_j(e,h).
\end{equation}
For a prescribed small $e_{cs,0}\in(I-\Pt u)X$, define
$\Phi_u(e,h)=\mathcal B(e_{cs,0},F_j(e,h))$.
We use the product norm
\[
 \norm{(e,h)}_{\mathcal Z_j}
       =\norm e_{\mathcal S_\nu}
                    +\delta_j^{-1/2}\norm h_{\mathcal Y_\nu}.
\]
These are complete spaces. In particular a locally uniform $X$ limit
and an $L^3L^{18/5}$ limit of its gradients agree as distributions;
the conjugate-paired real optical subspace is closed.
The density weight makes both directions of the feedback small.
Writing $E=\norm e_{\mathcal S_\nu}$ and
$H=\delta_j^{-1/2}\norm h_{\mathcal Y_\nu}$, the preceding estimates give
some $\eta_j\to0$ such that
\begin{align}
 \norm{\Phi_u(e,h)}_{\mathcal S_\nu}
 &\le C\bigl[\norm{e_{cs,0}}_X+\eta_jE+E^2+E^3
                         +\sqrt{\delta_j}H(1+E)\bigr],\label{eq:joint-u}\\
 \delta_j^{-1/2}\norm{\Phi_\rho(e)}_{\mathcal Y_\nu}
 &\le C\sqrt{\delta_j}(1+\lambda_0+E+E^2+E^3+E^4).
 \label{eq:joint-rho}
\end{align}
On interval $[k,k+1]$, $a_I(e)\le Ee^{-\nu k}$ and
$\norm h_{CH^1}\le\sqrt{\delta_j}He^{-\nu k}$.
Multiplying the forcing estimate by $e^{\nu k}$ gives the displayed
polynomials, including the linear feedback term. The frequency part
$(\omega_j-1)e$ is estimated directly in $X$.

On a ball of radius $r<1$ in $\mathcal Z_j$, the difference estimates
give Lipschitz constant at most
\[
 C\bigl[\eta_j+r+r^2+\sqrt{\delta_j}(1+r)^3\bigr].
\]
First choose $r$ small and then $j$ sufficiently large to make this less
than $1/2$. Fix a sufficiently small $\lambda_0>0$ and choose
$\norm{e_{cs,0}}_X$ small. Increasing $j$ if needed makes the constant
terms in~\eqref{eq:joint-u}--\eqref{eq:joint-rho} sum to at most $r/2$.
The map $(\Phi_u,\Phi_\rho)$ is therefore a contraction of the closed ball.
It has a fixed point satisfying both exact equations, with
\begin{equation}\label{eq:decay-perturbation}
 \norm{e(s)}_X+\norm{h(s)}_{H^1}\longrightarrow0
                       \qquad(s\to\infty).
\end{equation}
The spectral constants and projections in this argument belong to the one
fixed scalar profile. Only the bounded coefficient differences in
\eqref{eq:coefficient-convergence} are perturbed.

\subsection{Choosing compact initial data}

Let $\chi\in C_c^\infty(\R^3)$ equal one on a ball, and
$\chi_L(y)=\chi(y/L)$. For the selected $j$, the tail
$w_L=(\chi_L-1)Q_j$ tends to zero in $X$.
To verify this in the homogeneous norm, partition the tail into dyadic
annuli of radii $R\ge L$. Each annular term $f_R$ obeys
\[
 \norm{f_R}_2\le C R^{1/2},\qquad
 \norm{\nabla f_R}_2\le C R^{-1/2},\qquad
 \norm{f_R}_{\dH^\sigma}
 \le\norm{f_R}_2^{1-\sigma}\norm{\nabla f_R}_2^\sigma
 \le C R^{1/2-\sigma}.
\]
The two geometric series converge and tend to zero as $L\to\infty$,
because $\sigma>1/2$. Local finiteness identifies their sum with the
actual tail as a distribution.

Choose the central-stable datum
$e_{cs,0}=(I-\Pt u)w_L$. It is small for large $L$.
Lemma~\ref{lem:green} supplies an initial correction
$v_u\in\Ran\Pt u\subset C_c^\infty$, so
\begin{equation}\label{eq:compact-perturbed}
 Q_j+e(0)=\chi_LQ_j+v_u-\Pt u w_L\in C_c^\infty.
\end{equation}
The resulting physical data are
\begin{equation}\label{eq:physical-data}
 u_0(x)=\lambda_0^{-1}[Q_j+e(0)](x/\lambda_0),\qquad
 \rho_0(x)=\lambda_0^{-2}[R_j+h(0)](x/\lambda_0)=0.
\end{equation}
The optical datum is compact and smooth, despite the nonintegrable
mass tail of the profile before cancellation.

\subsection{Identification with the physical Cauchy solution}
\label{subsec:local-theory}

For completeness we give the local theory used to identify the constructed
solution. The physical equations at the selected parameters are
\begin{equation}\label{eq:special-system}
 i(\partial_t+\partial_z)u+\Delta u
       =-(1+i\delta_j)|u|^2u+\rho u,
 \qquad \rho_t=\delta_j|u|^4.
\end{equation}
The free optical group is
$S(t)f(x)=(e^{it\Delta}f)(x-te_z)$.
Its spatial norms equal those of the free Schr\"odinger group.
The free retarded Strichartz estimates~\cite[Theorem 2.3]{Tao}, with
derivatives, control
\[
 A_I(u)=\norm u_{C(I;H^1)}+
             \norm{\nabla u}_{L^3(I;L^{18/5})}
\]
from $H^1$ data and forcings in $L^{4/3}W^{1,3/2}+L^2W^{1,6/5}$.
The pairs $(4,3)$ and $(2,6)$ are admissible in dimension three.
The source normalization $\Delta/2$ changes to $\Delta$ by $t\mapsto2t$,
which changes the time norms only by fixed factors.
Estimate~\eqref{eq:mixed-GN} gives
$\norm u_{L^4(I;L^\infty)}\le C A_I(u)$.

For $|I|\le1$, H\"older and the product rule give
\begin{align}
 \norm{|u|^4}_{L^1H^1}
 &\le |I|^{3/4}\norm u_{CL^6}^3\norm u_{L^4L^\infty}
       +4|I|^{1/4}\norm u_{L^4L^\infty}^3\norm{\nabla u}_{CL^2},
       \label{eq:physical-quartic}\\
 \norm{\rho u}_{L^{4/3}W^{1,3/2}}
 &\le C|I|^{3/4}\norm\rho_{CH^1}\norm u_{CH^1},
       \label{eq:physical-potential}\\
 \norm{|u|^2u}_{L^2W^{1,6/5}}
 &\le C|I|^{1/2}\norm u_{CH^1}^3.
       \label{eq:physical-cubic}
\end{align}
All time norms here are on $I$. For the derivative in
\eqref{eq:physical-potential}, use
$\nabla\rho\in L^2,u\in L^6$ and
$\rho\in L^6,\nabla u\in L^2$; the undifferentiated product uses
$\rho\in L^2,u\in L^6$.
For the cubic derivative use two $L^6$ factors and one $L^2$ derivative;
for the undifferentiated term use two $L^6$ factors and one $L^2$ factor.
The quartic bounds are the same degree-four estimates as above.

Differences are estimated by factoring products. In particular, a
differentiated quartic difference is bounded by terms of the form
\[
 (|u|+|v|)^3|\nabla(u-v)|
 \quad\hbox{and}\quad
 (|u|+|v|)^2|u-v|(|\nabla u|+|\nabla v|).
\]
Three $L^4_tL^\infty_x$ factors and one $C_tL^2_x$ derivative give
$|I|^{1/4}$. The undifferentiated difference uses three $C_tL^6_x$
factors and one $L^4_tL^\infty_x$ difference, giving $|I|^{3/4}$.
Thus the optical Duhamel map and
$\rho=\rho_0+\delta_j\int|u|^4$ contract on a sufficiently short interval
in the class $A_I\times C(I;H^1)$, for each fixed $H^1$ datum.
This proves local existence and uniqueness in that class.

The reconstruction~\eqref{eq:ansatz} initially gives
$u\in C_{\rm loc}([0,T);X)$,
$\rho\in C_{\rm loc}([0,T);H^1)$ and the local gradient Strichartz
bound. On a compact physical interval the similarity time interval is
compact and all scale factors are bounded above and below.
In particular
\[
 \norm{\lambda^{-1}f(\cdot/\lambda)}_{\dH^a}
       =\lambda^{1/2-a}\norm f_{\dH^a}\quad(a=\sigma,1),
\]
whereas the density $L^2$ and $\dH^1$ factors are
$\lambda^{-1/2}$ and $\lambda^{-3/2}$.
Translations, phases and dilations are strongly continuous on these
spaces, as follows first on Schwartz functions and then by completion.
The background gradient belongs to $L^{18/5}$ by its uniform decay.
The exact coordinate calculation proves~\eqref{eq:special-system}
distributionally for this reconstruction.

To recover the missing mass norm, observe that $u\in CL^6$ and
$\rho\in CL^3$ on compact intervals, so the optical forcing belongs to
$CL^2$:
\[
 \norm{|u|^2u}_2=\norm u_6^3,
 \qquad \norm{\rho u}_2\le\norm\rho_3\norm u_6.
\]
Since~\eqref{eq:physical-data} belongs to $L^2$, its physical Duhamel
formula is in $CL^2$. It equals the reconstructed distributional solution:
their difference solves the free equation with zero forcing and zero
initial trace; spatial Fourier transformation and the Schr\"odinger
integrating factor, together with the transport translation, make it zero.
Consequently $u\in CH^1$ on every compact interval before $T$, and the
local uniqueness just proved applies.

The original density is therefore
\begin{equation}\label{eq:physical-positivity}
 \rho(t,x)=\delta_j\int_0^t |u(\tau,x)|^4\,d\tau\ge0
\end{equation}
almost everywhere. To identify the $H^1$ integral with this scalar
integral, use~\eqref{eq:physical-quartic} and spatial H\"older on bounded
boxes: the absolute spacetime integral is finite there. Fubini--Tonelli
then applies. Positivity concerns the actual density even though $h$ can
have either sign.

\subsection{Smooth persistence and divergence}

For each integer $m\ge3$, $H^m$ embeds into $L^\infty$ and is a product
algebra. The physical Duhamel and density integral maps are locally
Lipschitz on $C(I;H^m\times H^m)$, with a lifespan bounded below in
terms of the initial norms. These solutions agree with the preceding
$H^1$ solution by uniqueness.
On every $[0,T']$ with $T'<T$, the construction gives
$u\in L^4_tL^\infty_x$. Equation~\eqref{eq:physical-positivity} implies
\[
 \sup_{t\le T'}\norm{\rho(t)}_\infty
       \le\delta_j\int_0^{T'}\norm{u(t)}_\infty^4\,dt<\infty.
\]
The tame Sobolev product estimate~\cite[Appendix A, Lemma A.8]{Tao}
applied to~\eqref{eq:special-system} yields, for
$Y_m=\norm u_{H^m}+\norm\rho_{H^m}$,
\begin{equation}\label{eq:persistence}
 Y_m(t)\le Y_m(0)+C_{m,\delta_j}\int_0^t
 \bigl[1+\norm u_\infty^2+\norm u_\infty^3+\norm\rho_\infty\bigr]
                    Y_m(\tau)\,d\tau.
\end{equation}
The optical potential contributes
$\norm\rho_\infty\norm u_{H^m}+
\norm u_\infty\norm\rho_{H^m}$, the cubic contributes
$C\norm u_\infty^2\norm u_{H^m}$, and the quartic density source
contributes $C\norm u_\infty^3\norm u_{H^m}$.
The bound follows directly from the Duhamel norm inequality and does
not differentiate a norm through the unbounded Laplacian.
Its coefficient is integrable on $[0,T']$. The integral Gronwall
inequality gives a uniform $H^m$ bound there.
If a smooth lifespan ended at $S<T$, the constructed solution would
supply the same integrable coefficient beyond $S$. The uniform norm
bound and the norm-dependent local lifespan would extend it, a
contradiction. Thus every fixed $m$ persists on $[0,T)$; uniqueness
identifies all these solutions. Density nonnegativity is now pointwise.

\begin{proof}[Proof of Theorem~\ref{thm:main}]
Select the finite $j$, the scale $\lambda_0$ and the cutoff $L$ as above.
The data~\eqref{eq:physical-data} are Schwartz, with $\rho_0=0$.
We have obtained the original smooth solution on $[0,T)$, with all the
claimed Sobolev regularity and nonnegative density.
The norm $\norm{\nabla Q_j}_2$ is positive: a zero gradient would make
$Q_j$ constant, contrary to its nonzero value and decay.
By~\eqref{eq:decay-perturbation},
\[
 \norm{\nabla(Q_j+e(s))}_2\longrightarrow\norm{\nabla Q_j}_2>0.
\]
Undoing the dilation, phase and translation gives
\begin{equation}\label{eq:blowup}
 \norm{\nabla u(t(s))}_2
       =\lambda(s)^{-1/2}\norm{\nabla(Q_j+e(s))}_2
       \longrightarrow\infty,\qquad t(s)\uparrow T<\infty.
\end{equation}
The datum $u_0$ cannot vanish: uniqueness would then give the zero
optical solution, contradicting~\eqref{eq:decay-perturbation} and
$Q_j\ne0$.
The chosen coefficient is the strictly positive $\delta_j$, rather than
its zero limit. Also $\eps=c=c_g=1$, $K=2$, and $\alpha_5=0$ are
admissible in the original system. The Laplacian has all three elliptic
spatial directions, and the density translation was retained exactly.
A global $C_t(H^3\times H^3)$ solution would coincide with this one
before $T$ and have bounded gradient norm near $T$, contradicting
\eqref{eq:blowup}. This proves the asserted counterexample.
\end{proof}

\appendix
\section{Finite exact coefficient inequalities}\label{app:arithmetic}

Only a small additional arithmetic calculation is needed beyond the
quantitative scalar theorems in~\cite{DS}. Recall the constants from
Section~\ref{subsec:scalar-construction}:
\[
 \beta_*=\frac{772201763088846}{841768781900003},\qquad
 r_*=1.2\cdot10^{-6},\qquad
 \beta_{\min}=\beta_*-10^{-10},\qquad
 \beta\in[\beta_{\min},\beta_*+10^{-10}].
\]
For its published polynomial
$P_* =\sum_{n=0}^{50}c_nT_n$, set
\[
 U=\sum c_n,\qquad V=\sum(-1)^nc_n,\qquad W=\sum n^2c_n,
 \qquad |z|_1=|\Rea z|+|\Ima z|,
 \qquad \phi(\beta)=\frac{2i\beta}{1-i\beta}.
\]
The identities $T_n(1)=1$, $T_n(-1)=(-1)^n$ and $T_n'(1)=n^2$
follow from the Chebyshev recurrence and its derivative.
Also $|T_n(y)|\le1$ on $[-1,1]$, and
$|\phi|\le2$, $|\phi'|\le2$ for positive real $\beta$.
Consequently~\eqref{eq:profile-envelopes} follows from the three exact
rational comparisons
\begin{align}
 2\sum_{n=0}^{50}|c_n|_1+2\sum_{n=0}^{50}n^2|c_n|_1+r_*&<50,
 \label{eq:finite-envelope}\\
 \Rea[V-U+\phi(\beta_*)W]+2\cdot10^{-10}|W|_1+r_*&<-1,
 \label{eq:finite-origin}\\
 \frac{4\beta_{\min}^2}{1+\beta_{\min}^2}|W|^2&>\frac1{100}.
 \label{eq:finite-tail}
\end{align}
The last bound uses that $\beta^2/(1+\beta^2)$ increases for $\beta>0$.
In particular $|F(1)|>1/10-r_*>0$.

The companion package supplies the 51 complex rational coefficients from
Appendix~A of the pinned version of~\cite{DS}, and the standalone script
\texttt{code/check\_constants.py}. Its arithmetic uses Python's exact
\texttt{Fraction} type. It checks~\eqref{eq:finite-envelope}--\eqref{eq:finite-tail},
the self-map inequality,~\eqref{eq:profile-contraction}, the two obstruction
margins and the positive-$\delta$ bounds in~\eqref{eq:P-bound}--\eqref{eq:P-difference}.
The endpoint sums are checked independently by constructing the
Chebyshev polynomials in the monomial basis. Input hashes and a replay
command are in the supplement's \texttt{code/README.md}. No floating-point comparison,
spatial discretization or PDE simulation is involved.
This finite calculation does not reproduce the full computer-assisted
inverse and residual proof of~\cite{DS}; those published theorems remain
external dependencies. The bound $10^{-23}$ used for the profile
perturbation is not a numerical threshold for the full Cauchy
construction, which also requires the qualitative choices after
\eqref{eq:joint-rho}.

\section*{AI use}
AI assistance was used for mathematical exploration, preparation of the
proof and manuscript, exact-arithmetic code, and internal proof checking.
Two further AI reviews in fresh contexts, a reader simulation and a
referee simulation, examined the same anonymized draft. They shared the
working environment. These activities do not constitute independent
human verification or formal proof certification. The argument uses the published scalar
computer-assisted construction of Donninger and Sch\"orkhuber and the
linear estimates of Li as cited above.

\begin{thebibliography}{9}
\bibitem{DLS}
\'E. Dumas, D. Lannes and J. Szeftel,
\emph{Variants of the Focusing NLS Equation: Derivation, Justification,
and Open Problems Related to Filamentation},
in \emph{Laser Filamentation}, CRM Series in Mathematical Physics,
Springer, 2016, pp.~19--75.
\href{https://doi.org/10.1007/978-3-319-23084-9_2}{doi:10.1007/978-3-319-23084-9\_2}.
\bibitem{DS}
R. Donninger and B. Sch\"orkhuber,
\emph{Self-Similar Blowup for the Cubic Schr\"odinger Equation},
arXiv:2406.16597v3, 9 December 2025.
\url{https://arxiv.org/abs/2406.16597v3}.
\bibitem{Li}
Z. Li,
\emph{On Stability of Self-Similar Blowup for Mass Supercritical NLS},
arXiv:2304.02078v2, 27 December 2024.
\url{https://arxiv.org/abs/2304.02078v2}.
\bibitem{Tao}
T. Tao,
\emph{Nonlinear dispersive equations: local and global analysis},
author-hosted text.
\url{https://www.math.ucla.edu/~tao/preprints/chapter.pdf}.
\end{thebibliography}
\end{document}
